\chapter{Forward Kinematics}
\label{ch:fk}

\section{2R}
\label{sec:2R}

% =====================================================================
%  Forward kinematics: a more generalizable approach
%  Goes into chapters/03_forward_kinematics.tex
%  Requires: tikz libraries calc, angles, quotes (already in preamble)
% =====================================================================

% --- Sketch helpers --------------------------------------------------
% Coordinate frame (drawn in colour so that an axis lying along a link
% stays visible): y-axis (horizontal) and z-axis (vertical),
% rotated as a whole with  \pic[rotate=<angle>] at (<point>) {axes};
\tikzset{
  axes/.pic={
    \draw[-{Stealth[length=4pt]}, semithick, blue!70!black]
      (0,0) -- (0.7,0);
    \draw[-{Stealth[length=4pt]}, semithick, blue!70!black]
      (0,0) -- (0,0.7);
  },
  fksketch/.style={scale=1.1, thick, baseline=(current bounding box.center)},
}
% Geometry used in all sketches (illustrative values only):
% R_1 = 2, R_2 = 1.6, varphi_1 = 35 deg, varphi_2 = 40 deg.
\newcommand{\fkpoints}{%
  \coordinate (O) at (0,0);
  \coordinate (J) at (125:2);                 % joint 2, end of link 1
  \coordinate (E) at ($(J)+(165:1.6)$);       % end of link 2
}

% Variable definitions added during transcription -- please check.
The arm is treated in the plane of its two links, with $y$ the horizontal
and $z$ the vertical coordinate. $R_1$ and $R_2$ are the lengths of the
first and second link, $\varphi_1$ and $\varphi_2$ are the joint angles,
and positions are written as $(1, y, z)\T$, so that translations can be
expressed as matrix products as well. The primes on $\vect{r}_2$ indicate
the basis in which it is expressed. 

\begin{center}
\begin{tikzpicture}[fksketch]
  \fkpoints
  \draw (-3.2,0) -- (1,0);
  \draw (O) -- (J) -- (E);
  \node at (E) {$\times$};
  \pic at (O) {axes};
  \pic[rotate=35] at (J) {axes};
\end{tikzpicture}
\end{center}

First determine the position $\vect{r}_2$ in the basis attached to the
second link:

\noindent
\begin{minipage}{0.35\textwidth}
  \centering
  \begin{tikzpicture}[fksketch]
    \fkpoints
    \draw (J) -- (E);
    \node at (E) {$\times$};
    \pic[rotate=75] at (J) {axes};
  \end{tikzpicture}
\end{minipage}\hfill
\begin{minipage}{0.6\textwidth}
  \[
    \vect{r}_2^{\prime\prime\prime} =
    \begin{pmatrix} 1 \\ 0 \\ R_2 \end{pmatrix}
  \]
\end{minipage}

\medskip

Change basis to the basis where the z axis aligns with the next link. 

\noindent
\begin{minipage}{0.35\textwidth}
  \centering
  \begin{tikzpicture}[fksketch]
    \fkpoints
    \draw[dashed, thin] (J) -- ($(J)+(125:1)$) coordinate (Z2);
    \draw (J) -- (E);
    \node at (E) {$\times$};
    \pic[rotate=35] at (J) {axes};
    \pic[draw, thin, "$\varphi_2$", angle radius=14pt,
         angle eccentricity=1.7] {angle = Z2--J--E};
  \end{tikzpicture}
\end{minipage}\hfill
\begin{minipage}{0.6\textwidth}
  \[
    \vect{r}_2^{\prime\prime} =
    \begin{pmatrix}
      1 & 0 & 0 \\
      0 & \cos\varphi_2 & -\sin\varphi_2 \\
      0 & \sin\varphi_2 & \cos\varphi_2
    \end{pmatrix}
    \begin{pmatrix} 1 \\ 0 \\ R_2 \end{pmatrix}
    =
    \begin{pmatrix} 1 \\ -R_2\sin\varphi_2 \\ R_2\cos\varphi_2 \end{pmatrix}
  \]
\end{minipage}

\medskip
Change basis to the basis that is still aligned with the link but sits at the bottom joint.
This makes the base change a mere translation.

\noindent
\begin{minipage}{0.35\textwidth}
  \centering
  \begin{tikzpicture}[fksketch]
    \fkpoints
    \draw (-3.2,0) -- (1,0);
    \draw (O) -- (J) -- (E);
    \fill (J) circle (1.5pt);
    \node at (E) {$\times$};
    \pic[rotate=35] at (O) {axes};
  \end{tikzpicture}
\end{minipage}\hfill
\begin{minipage}{0.6\textwidth}
  \[
    \vect{r}_2^{\prime} =
    \begin{pmatrix}
      1   & 0 & 0 \\
      0   & 1 & 0 \\
      R_1 & 0 & 1
    \end{pmatrix}
    \vect{r}_2^{\prime\prime}
    =
    \begin{pmatrix} 1 \\ -R_2\sin\varphi_2 \\ R_2\cos\varphi_2 + R_1 \end{pmatrix}
  \]
\end{minipage}

\medskip
Change basis to standard y-z basis with origin at base joint.

\noindent
\begin{minipage}{0.35\textwidth}
  \centering
  \begin{tikzpicture}[fksketch]
    \fkpoints
    \draw (-3.2,0) -- (1,0);
    \coordinate (Z0) at (0,1);
    \draw (O) -- (J) -- (E);
    \node at (E) {$\times$};
    \pic at (O) {axes};
    \pic[draw, thin, "$\varphi_1$", angle radius=16pt,
         angle eccentricity=1.6] {angle = Z0--O--J};
  \end{tikzpicture}
\end{minipage}\hfill
\begin{minipage}{0.6\textwidth}
  \[
    \vect{r}_2 =
    \begin{pmatrix}
      1 & 0 & 0 \\
      0 & \cos\varphi_1 & -\sin\varphi_1 \\
      0 & \sin\varphi_1 & \cos\varphi_1
    \end{pmatrix}
    \vect{r}_2^{\prime}
  \]
\end{minipage}

\begin{align*}
  \vect{r}_2
  &= \begin{pmatrix}
       1 \\
       -R_2\cos\varphi_1\sin\varphi_2 - R_2\cos\varphi_2\sin\varphi_1
         - R_1\sin\varphi_1 \\
       -R_2\sin\varphi_1\sin\varphi_2 + R_2\cos\varphi_1\cos\varphi_2
         + R_1\cos\varphi_1
     \end{pmatrix} \\
  &= \begin{pmatrix}
       1 \\
       -R_1\sin\varphi_1 - R_2\sin(\varphi_1 + \varphi_2) \\
       R_1\cos\varphi_1 + R_2\cos(\varphi_1 + \varphi_2)
     \end{pmatrix}
\end{align*}

The entire transformation:
\[
  \vect{r}_2 =
  \begin{pmatrix}
    1 & 0 & 0 \\
    0 & \cos\varphi_1 & -\sin\varphi_1 \\
    0 & \sin\varphi_1 & \cos\varphi_1
  \end{pmatrix}
  \begin{pmatrix}
    1   & 0 & 0 \\
    0   & 1 & 0 \\
    R_1 & 0 & 1
  \end{pmatrix}
  \begin{pmatrix}
    1 & 0 & 0 \\
    0 & \cos\varphi_2 & -\sin\varphi_2 \\
    0 & \sin\varphi_2 & \cos\varphi_2
  \end{pmatrix}
  \vect{r}_2^{\prime\prime\prime}
\]

% =====================================================================
%  chapters/03_forward_kinematics.tex -- 4R arm
%  Goes after the section "A more generalizable approach".
%  Requires: tikz (calc, angles, quotes, arrows.meta), subcaption
% =====================================================================

\section{Forward kinematics: 4R arm}
\label{sec:4r-fk}

The arm consists of a base that rotates about the vertical $z$-axis by
the angle $\varphi_0$, followed by three links of lengths $R_1$, $R_2$
and $R_3$ with joint angles $\varphi_1$, $\varphi_2$ and $\varphi_3$
acting in a vertical plane. The first joint sits at height $h$ above the
ground, and $D_0$, $D_1$ and $D_2$ are fixed lateral offsets between the
joints. All joint angles are signed, with counterclockwise positive;
angles drawn in the sketches are unsigned magnitudes. Positions are
written in homogeneous coordinates $(1, x, y, z)\T$.


\Cref{fig:4r-structure} shows the assumed general structure, with
$\varphi_0 = 0$ in the top view.

\begin{figure}[tbp]
  \centering
  \begin{subfigure}[b]{0.5\textwidth}
    \centering
    \begin{tikzpicture}[thick, scale=1.2]
      \draw[-{Stealth[length=4pt]}, semithick] (-1.2,2.6) -- (-0.7,2.6) node[right] {$x$};
      \draw[-{Stealth[length=4pt]}, semithick] (-1.2,2.6) -- (-1.2,3.1) node[above] {$z$};
      \draw (-1,0) -- (3.4,0);
      \draw (-0.4,0) rectangle (0.4,0.4);
      \draw[|<->|, thin] (-0.7,0) -- node[left] {$h$} (-0.7,0.4);
      \coordinate (S) at (0,0.4);
      \coordinate (E) at ($(S)+(65:2)$);
      \coordinate (W) at ($(E)+(15:1.6)$);
      \coordinate (T) at ($(W)+(-40:0.9)$);
      \coordinate (U) at ($(S)+(0,1)$);
      \coordinate (X1) at ($(E)+(65:0.8)$);
      \coordinate (X2) at ($(W)+(15:0.7)$);
      \draw[dashed, thin] (S) -- (U);
      \draw[dashed, thin] (E) -- (X1);
      \draw[dashed, thin] (W) -- (X2);
      \draw (S) -- node[right, pos=0.55] {$R_1$} (E)
                -- node[below] {$R_2$} (W)
                -- node[below left] {$R_3$} (T);
      \foreach \p in {S,E,W} \fill (\p) circle (1.3pt);
      \node at (T) {$\times$};
      \pic[draw, thin, "$\varphi_1$", angle radius=16pt, angle eccentricity=1.5]
        {angle = E--S--U};
      \pic[draw, thin, "$\varphi_2$", angle radius=14pt, angle eccentricity=1.6]
        {angle = W--E--X1};
      \pic[draw, thin, "$\varphi_3$", angle radius=12pt, angle eccentricity=1.7]
        {angle = T--W--X2};
    \end{tikzpicture}
    \caption{Side view.}
  \end{subfigure}\hfill
  \begin{subfigure}[b]{0.46\textwidth}
    \centering
    \begin{tikzpicture}[thick, scale=1.1]
      \draw[-{Stealth[length=4pt]}, semithick] (-0.6,1.6) -- (-0.1,1.6) node[right] {$x$};
      \draw[-{Stealth[length=4pt]}, semithick] (-0.6,1.6) -- (-0.6,2.1) node[above] {$y$};
      \draw[thin] (0,0) circle (0.3);
      \fill (0,0) circle (1.3pt);
      \draw[dashed, thin] (0,0) -- (3.6,0);
      \draw (0,0) -- (0,0.35) -- (1,0.35) -- (1,0.7) -- (2,0.7) -- (2,1.05) -- (2.9,1.05);
      \node at (2.9,1.05) {$\times$};
      \draw[|<->|, thin] (-0.45,0) -- node[left] {$D_0$} (-0.45,0.35);
      \draw[|<->|, thin] (0.75,0.35) -- node[left] {$D_1$} (0.75,0.7);
      \draw[|<->|, thin] (1.75,0.7) -- node[left] {$D_2$} (1.75,1.05);
      \draw[dotted, thin] (0,0) -- (2.9,1.05) coordinate (Tt);
      \coordinate (O) at (0,0); \coordinate (Xa) at (3.6,0);
      \pic[draw, thin, "$\delta\varphi_0$", angle radius=70pt, angle eccentricity=1.3]
        {angle = Xa--O--Tt};
    \end{tikzpicture}
    \caption{Top view, $\varphi_0 = 0$.}
  \end{subfigure}
  \caption{Assumed general structure of the 4R arm.}
  \label{fig:4r-structure}
\end{figure}

Again we work our way slowly with basis transformations, from the trivial
description of the position $\vect{r}'$ of the end effector to the
description $\vect{r}$ in the general basis:
\begin{align*}
  \vect{r} =\;
  &\begin{pmatrix}
     1 & 0 & 0 & 0 \\
     0 & \cos\varphi_0 & -\sin\varphi_0 & 0 \\
     0 & \sin\varphi_0 & \cos\varphi_0 & 0 \\
     0 & 0 & 0 & 1
   \end{pmatrix}
   \begin{pmatrix}
     1 & 0 & 0 & 0 \\
     0 & \cos\varphi_1 & 0 & -\sin\varphi_1 \\
     D_0 & 0 & 1 & 0 \\
     h & \sin\varphi_1 & 0 & \cos\varphi_1
   \end{pmatrix} \\
  &\begin{pmatrix}
     1 & 0 & 0 & 0 \\
     0 & \cos\varphi_2 & 0 & -\sin\varphi_2 \\
     D_1 & 0 & 1 & 0 \\
     R_1 & \sin\varphi_2 & 0 & \cos\varphi_2
   \end{pmatrix}
   \begin{pmatrix}
     1 & 0 & 0 & 0 \\
     0 & \cos\varphi_3 & 0 & -\sin\varphi_3 \\
     D_2 & 0 & 1 & 0 \\
     R_2 & \sin\varphi_3 & 0 & \cos\varphi_3
   \end{pmatrix}
   \vect{r}'
\end{align*}
with
\[
  \vect{r}' = \begin{pmatrix} 1 \\ 0 \\ 0 \\ R_3 \end{pmatrix} .
\]
